\section{Grants}\label{sec:grants}

\subsection{The model}

A grant over an action type $\alpha$ is a list of caveats, each a predicate on
$\alpha$. Figure~\ref{fig:model} is the readable Lean specification
(\file{formal/AttenuatingGrants/Grant.lean}), quoted from the repository. An
action is allowed when every caveat holds, so the empty grant allows every
action. A batch of actions, which is how every context of a Soroban
authorization tree is checked, is allowed when every action in it is. One
grant is narrower than another when every action the first allows, the second
allows too, and attenuation adds caveats.

\begin{figure}[t]
\begin{lstlisting}[language=Lean]
structure Grant (α : Type) where
  caveats : List (α → Bool)

def Grant.allows (g : Grant α) (a : α) : Bool :=
  g.caveats.all (· a)

def Grant.allowsAll (g : Grant α) (actions : List α) : Bool :=
  actions.all g.allows

def Grant.narrower (g h : Grant α) : Prop :=
  ∀ a, g.allows a = true → h.allows a = true

def Grant.attenuate (g : Grant α) (c : α → Bool) : Grant α :=
  { caveats := c :: g.caveats }
\end{lstlisting}
\caption{The readable grant model.}\label{fig:model}
\Description{Lean definitions of Grant, allows, allowsAll, narrower and attenuate.}
\end{figure}

\paragraph{Two documents.} The grant signer and the agent host split the
operations on grants between two interfaces. The \emph{admin document}
issues a root grant, revokes a grant (and with it every grant delegated from
it), rotates a holder's key and refills a budget. The \emph{agent interface}
acts under a grant and delegates: \code{delegate(add)} creates a child grant
whose caveats are the caller's followed by \code{add}, so the child allows
only what the caller allows, whatever \code{add} says. No agent operation
reads or edits a grant. At the relay, grants are a configuration file that
the admin side writes; agents only send requests.

\subsection{Four vocabularies}

Four action vocabularies plug into the model (Table~\ref{tab:vocab}). Each
defines an action type, a caveat type and a denotation from caveats to
predicates. Identifiers (contract ids, function names, wasm hashes, hosts,
key ids, network ids, path components) are byte strings compared by exact
equality. The evaluator normalizes nothing: lower-casing hosts and resolving
\code{..}, symlinks and case happen in the host before an action exists
(\S\ref{sec:host}).

\begin{table}[t]
\caption{The four action vocabularies.}\label{tab:vocab}
\small
\begin{tabularx}{\linewidth}{@{}l>{\raggedright\arraybackslash}X>{\raggedright\arraybackslash}X@{}}
\toprule
& \textbf{Action} & \textbf{Caveats} \\
\midrule
Chain & \code{invoke} contract function args, or \code{createContract} wasmHash
  & \code{contractIn}, \code{functionIn}, \code{argI128AtMost} $i$ $n$, \code{argEq} $i$ $v$ \\
Files & path (a list of components), mode $\in\{$read, write, execute$\}$
  & \code{pathUnder} prefix, \code{modeIn} \\
Network & host, port, scheme $\in\{$http, https$\}$
  & \code{hostIn}, \code{portIn}, \code{schemeIn} \\
Signing & key, network, payload $\in\{$transaction, authEntry, message$\}$
  & \code{keyIn}, \code{networkIn}, \code{payloadIn} \\
\bottomrule
\end{tabularx}
\end{table}

\paragraph{Chain.} Arguments follow Soroban's values: \code{i128},
\code{u64}, \code{bool}, \code{bytes}, \code{address} and \code{symbol} are
modelled, and an argument of any other type is \code{other} carrying its full
XDR encoding, whose leading discriminant names its type. Values of different
constructors are never equal, so an \code{argEq} caveat holds only for an
argument of its value's constructor\claim{arg-kind}, and an unmodelled
argument never matches a \code{bytes} caveat, whatever its
encoding\claim{other-not-bytes}. Every chain caveat is false on
\code{createContract}, so any chain grant with at least one caveat denies
contract deployment\claim{create-contract-denied}.

\paragraph{Caveats apply to every context alike.} A chain grant is checked
against every context of an authorization tree, and each caveat reads the same
fields of each context: \code{argI128AtMost i n} and \code{argEq i v} read
argument $i$ of every context, whatever its contract and
function\claim{positional-caveats}. Two consequences follow. A grant with
\code{contractIn [R, T]} and \code{functionIn [swap, transfer]} allows the
cross product, \code{R.transfer} and \code{T.swap} included. And a tree whose
calls carry amounts at different positions cannot have each amount capped: a
router's \code{forward} takes its amount as argument 3 and the token
\code{transfer} it nests as argument 2, so one \code{argI128AtMost} caps the
wrong argument of one of them. The evaluation's router workload therefore
bounds amounts by its allowance alone (\S\ref{sec:eval}). A caveat that
applies only to contexts of one contract and function would close this; the
generic model takes any predicate, so Theorems~\ref{thm:atten}
and~\ref{thm:chain} would be unchanged, but the vocabulary does not have one.

\paragraph{Files.} Paths are lists of components, so a grant under
\code{/home/a} never covers \code{/home/ab}\claim{path-components}. A
read-only grant is \code{modeIn [read]}.

\paragraph{Network and signing.} A network action is the host name the agent
requested, its port and its scheme; a signing action is a key id, a network
passphrase and a payload kind. Each caveat is membership of one field in a
list.

\subsection{Budgets}\label{sec:budgets}

A grant may carry an allowance of one asset. Budgets are enforced at the
Soroban boundary, where the grant signer meters transfers of one Stellar Asset
Contract per root grant; the agent host and the relay do not meter.
Allowances live on the delegation tree: the ledger holds one node per grant with a parent link, an
allowance and \code{spent}, the total of every successful spend at or under
the node. A spend at node $i$ walks from $i$ to its root and fails at the
first node whose \code{spent} plus the amount would exceed its allowance;
otherwise it charges every node on the walk. Either the spend succeeds and is
charged to the spender and every ancestor, or it fails and changes nothing.

Allowances may oversubscribe: children's allowances may sum to more than their
parent's, and delegation accepts any allowance. Because every ancestor is
charged in the same step, spending under a node still never exceeds that
node's allowance. The model's walk stops at the first node, counting from the
spender, that cannot take the amount, and a failed spend names it; that node
is where a refill request goes, up the delegation chain to the human who holds
the admin document. The proven statement about the name is weaker: it is a
node on the spender's chain that cannot take the amount (Theorem~\ref{thm:budget}).
Refilling and issuing a new root are admin operations; an agent can only
append a child node by delegating. Allowances are cumulative and never reset.

Figure~\ref{fig:budget} shows two runs from a root of 10 with two children of
8 each and nothing spent. In the first, after node~1 spends 8, node~2's spend
of 3 fails at the root, which has 2 left, and its spend of 2 succeeds. In the
second, node~1 delegates an allowance of 20, more than its own 8, to a new
node~3, whose spend of 9 is refused and whose spend of 8 succeeds, charging
node~3, node~1 and the root.

\begin{figure}[t]
\centering
\input{budget-tree}
\caption{Oversubscribed allowances. The \lean{example}s at the end of
\file{Budget.lean} check by evaluation the first refusal's error, and each
run's final ledger and log of successful spends; ``spent'' lists each node's
spent after the run, from node~0. A refused spend changes no
node.}\label{fig:budget}
\Description{A delegation tree: root with allowance 10, two children with
allowance 8 each, and a node with allowance 20 delegated under the first
child, with two runs of spends and their outcomes.}
\end{figure}
